\chapter{Estado de redacción y pendientes de revisión}
\label{ap:pendientes}

La etapa vigente es exploratoria y pensionaria: no se ha aprobado pregunta
definitiva, hipótesis principal, outcome, definición de vulnerabilidad ni modelo.
Los pendientes predictivos anteriores se conservan como posibilidades futuras,
no como obligaciones de esta etapa. La matriz de cumplimiento v1 distingue
trabajo provisional, parcial y comparaciones no viables. Primero se requiere
profundizar contratos semánticos de las columnas candidatas, número declarado
de pensiones, montos salariales mes/año, titularidad e instrumentos.

Los capítulos~\ref{cap:introduccion}--\ref{cap:datos} reúnen un borrador
de introducción, contexto pensionario, datos y auditorías. Las tablas se generan
desde agregados verificados y las figuras conservan estado de borrador editorial.
La revisión extensa de literatura y la homologación global siguen pendientes;
el capítulo~\ref{cap:resultados} contiene B01 revisado, el bloque parcial
independiente de K 2021 y la integración provisional con comienzo del EDA
conjunto, sin aceptación académica automática.

Antes de fijar la plantilla se deben confirmar las normas UDLAP con el
departamento y la dirección: portada, programa, márgenes, tipografía,
interlineado, numeración, referencias y entrega. La guía pública acepta
LaTeX, pero no sustituye esa confirmación \citep{udlap2024guia}.
Esta plantilla es provisional.

Faltan datos de autoría, dirección y jurado, título aprobado, resumen y revisión
bibliográfica completa. Permanecen abiertas población predictiva, definición
económica, precios, horizontes, anomalías, modalidades, atrición y ponderación.
Resultados y conclusiones requieren evidencia calculada y revisada. Cada
checkpoint mantendrá sincronizados análisis y manuscrito; disponer de un
borrador no cierra la aceptación metodológica.

El complemento pensionario conserva ocho conceptos propuestos, un catálogo
histórico selectivo y matrices íntegras. La integración provisional está
autorizada y ejecutada; para homologar todas las variables faltan validación
por pregunta, titular, filtros, unidades, periodos,
fuentes sin tipo y productos ausentes. Incorporar 2024 requiere aprobación
de alcance. Las pruebas de claves no resuelven esos pendientes ni aprueban
una muestra principal.

La auditoría específica de titularidad amplía el cotejo a las seis rondas,
sin cerrar controles tempranos ausentes ni discrepancias de punteros. Los
faltantes permanecen explícitos; la base completa permite continuar análisis
con contratos validados sin esperar al cierre de todos los bloques independientes.

La revisión B01 conserva abiertas la causa de las discrepancias de
\variable{n_res}, los filtros laborales y la elegibilidad parcial de hermanos
vivos de 2012. Las 91 comparaciones conjuntas y sus tablas están disponibles;
534 universos de indicadores reconstruidos no equivalen a validar todos los
filtros de preguntas originales. El bloque K 2021 describe 338 variables,
con salarios de periodicidad pendiente y universos brutos de viudez no
reconstruidos en los productos. Las demás preguntas y fuentes continúan
registradas como pendientes; no se omiten por falta de relación pensionaria.


\section*{Protección económica: escritura y validación pendientes}
El EDA de pensiones y donativos públicos de personas de 60 años o más se presenta
en las secciones~\ref{sec:proteccion-metodo} y~\ref{sec:proteccion-resultados}.
Su ejecución no cierra la clasificación: continúan controles históricos, punteros,
coincidencias y fuentes mixtas. Faltan ruralidad estricta, diseño/ponderación,
precios y equivalencia histórica de programas. Los borradores de figuras y tablas
requieren revisión editorial; no se aprobaron estimando, muestra predictiva ni
modelos. Los anexos computacionales conservan el detalle fuera de la selección
del manuscrito.
% Conciliación documental: no implica aprobación de nuevas decisiones.
\section{Recuperación del documento original y cobertura nacional}
La pregunta prospectiva se conserva como antecedente y línea futura opcional.
Solo para una eventual reactivación de esa línea faltaría conciliar cuartil frente a quintil,
la población y los cortes de referencia, los empates, precios y ponderación.
Las H1--H3 del original eran marcadores; las hipótesis de la propuesta anterior permanecen
como antecedentes sin aprobación actual. Los siete capítulos del índice original y las normas UDLAP
requieren revisión editorial. No se adopta automáticamente el panel de seis
entrevistas directas ni los umbrales de acumulación de deterioro.
El catálogo nacional documenta oferta de las 32 entidades mediante cortes
discontinuos; quedan pendientes reglas individuales, fechas de continuidad,
montos y vacíos históricos. No identifica receptores locales ni permite
resolver desconocidos por residencia, edad o importe.

\section{Conservación documental antes del retiro de históricos}
El informe general permite consultar las contribuciones de los 610 documentos
previos, con tablas, interpretaciones, límites y referencias de procedencia.
La evidencia documental se conserva con su fecha, universo y estado originales;
las propuestas predictivas anteriores no son decisiones vigentes.
Los once notebooks con contenido histórico único se archivan en investigación,
y las doce ejecuciones previas se mantienen como evidencia intermedia.
Las capturas de revisión del PDF se conservan localmente. El retiro físico de los históricos reemplazados fue autorizado y ejecutado
tras verificar correspondencias, respaldo y excepciones. Los nuevos análisis
justificados tendrán notebooks independientes en la misma carpeta y contribuirán
al único informe acumulativo, sin fusionarse automáticamente con los existentes.
